% !PRINTED 621
\weberpage{621}
%<<CONT
und hierin sind die Co\"efficienten
\[
f(1), f'(1), \tfrac12 f''(1), \cdots, \dfrac{1}{\Pi(\mu-1)} f^{(\mu-1)}(1)
\]
ganze rationale Zahlen, n\"amlich
\begin{align*}
\dfrac{1}{\Pi(\mu-1)} f^{(\mu-1)}(1) &= x_{\mu-1} \\[6pt]
\dfrac{1}{\Pi(\mu-2)} f^{(\mu-2)}(1) &= x_{\mu-2}+(\mu-1)x_{\mu-1} \\[6pt]
\dfrac{1}{\Pi(\mu-3)} f^{(\mu-3)}(1) &= x_{\mu-3}+(\mu-2)x_{\mu-2}+\dfrac{(\mu-1)(\mu-2)}{1.2}x_{\mu-1} \tag{5}\\[6pt]
\cdots\cdots\cdots &\cdots\cdots\cdots\cdots\cdots\cdots\cdots \\[6pt]
f(1) &= x_0+x_1+\cdots\cdots\cdots\cdots+x_{\mu-1}.
\end{align*}
Da nun $\omega$ durch $q$ und folglich durch $\sigma^\mu$ theilbar ist, so
muss $f(1)$ nach (4) durch $q$ theilbar sein. Da aber $\mu>1$ und
folglich $q:\sigma$ noch durch $\sigma$ theilbar ist, so folgt, dass auch
\[
f'(1) - \dfrac{\sigma}{1.2} f''(1) \ldots
\]
durch $\sigma$, folglich $f'(1)$ durch $q$ theilbar ist, und wenn man so
weiter schliesst, ergiebt sich, dass die s\"ammtlichen rationalen
Zahlen (5) durch $q$ theilbar sind, und folglich sind auch die
$x_{\mu-1}, x_{\mu-2}, \ldots, x_0$ durch $q$ theilbar. Damit ist der Satz II. vollst\"andig bewiesen.

Anstatt der Reihe (1) kann man auch eine beliebige Reihe von
$\mu$ auf einander folgenden Potenzen von $r$ als Basis von $\mathfrak{o}$ nehmen:
\weq{6}{r^a, r^{a+1}, r^{a+2}, \ldots, r^{a+\mu-1},}
die sich ergiebt, wenn man alle Glieder der Reihe (1) mit $r^a$
multiplicirt.

Die Richtigkeit dieser Bemerkung ergiebt sich daraus, dass
die $r^{\pm a}$ ganze Zahlen sind, und dass also $\omega$ und $r^a\omega$ stets gleichzeitig ganze oder gebrochene Zahlen sind.

Beispielsweise bilden die Zahlen
\weq{7}{r^{-\frac12\mu+1}, r^{-\frac12\mu+2}, \ldots, r^{-1}, 1, r, r^2, \ldots, r^{\frac12\mu}}
eine Basis von $\mathfrak{o}$.

Damit ist auch die Grundzahl $\Delta$ des K\"orpers $\Omega_m$ bestimmt.
Sie ist nach der allgemeinen Definition (\S.\,145) gleich der Discirminante % [sic] "Discirminante" as printed; "Discriminante" is meant
\[
\Delta(1, r, r^2, \ldots, r^{\mu-1}),
\]
d.\,h. gleich der in \S.\,169 bestimmten Discriminante der Kreistheilungsgleichung.
